\begingroup
\shorthandoff{:;!?}
\usetikzlibrary{calc}
\resizebox{\textwidth}{!}{%
\begin{tikzpicture}[
  font=\sffamily,
  x=1cm, y=1cm,
  badge/.style={
    circle, minimum size=1.18cm,
    draw=black!30, line width=0.55pt
  },
  labfleche/.style={
    font=\sffamily\tiny, text=black!55, align=center, inner sep=1.5pt
  },
  lien/.style={
    -{Stealth[length=2.1mm, width=1.55mm]},
    draw=black!65, line width=0.7pt,
    dash pattern=on 1.15pt off 1.45pt
  }
]

% Pictogrammes simples (impression noir et blanc)
\tikzset{
  pics/icoDossier/.style={code={
    \fill[black!18] (-0.22,-0.14) rectangle (0.24,0.08);
    \fill[black!18] (-0.22,0.08) -- (-0.08,0.20) -- (0.06,0.08);
    \draw[black!55, line width=0.45pt] (-0.22,-0.14) rectangle (0.24,0.08);
    \draw[black!55, line width=0.45pt] (-0.22,0.08) -- (-0.08,0.20) -- (0.06,0.08);
  }},
  pics/icoPage/.style={code={
    \fill[white] (-0.16,-0.22) rectangle (0.16,0.22);
    \draw[black!55, line width=0.45pt] (-0.16,-0.22) rectangle (0.16,0.22);
    \draw[black!45, line width=0.4pt] (-0.08,0.10) -- (0.08,0.10);
    \draw[black!45, line width=0.4pt] (-0.08,0.00) -- (0.08,0.00);
    \draw[black!45, line width=0.4pt] (-0.08,-0.10) -- (0.04,-0.10);
  }},
  pics/icoListe/.style={code={
    \foreach \y in {0.12, 0.00, -0.12} {
      \fill[black!45] (-0.16,\y) circle (1.15pt);
      \draw[black!50, line width=0.45pt] (-0.08,\y) -- (0.18,\y);
    }
  }},
  pics/icoArbre/.style={code={
    \draw[black!55, line width=0.5pt] (0,0.16) -- (0,0.02);
    \draw[black!55, line width=0.5pt] (0,0.02) -- (-0.16,-0.14);
    \draw[black!55, line width=0.5pt] (0,0.02) -- (0.16,-0.14);
    \fill[black!20] (0,0.18) circle (3.1pt);
    \fill[black!20] (-0.16,-0.16) circle (2.5pt);
    \fill[black!20] (0.16,-0.16) circle (2.5pt);
    \draw[black!55, line width=0.4pt] (0,0.18) circle (3.1pt);
    \draw[black!55, line width=0.4pt] (-0.16,-0.16) circle (2.5pt);
    \draw[black!55, line width=0.4pt] (0.16,-0.16) circle (2.5pt);
  }},
  pics/icoClef/.style={code={
    \draw[black!55, line width=0.7pt] (0,0.10) circle (4.2pt);
    \fill[white] (0,0.10) circle (1.6pt);
    \draw[black!55, line width=0.7pt] (0,0.02) -- (0,-0.20);
    \draw[black!55, line width=0.7pt] (0,-0.08) -- (0.08,-0.08);
    \draw[black!55, line width=0.7pt] (0,-0.16) -- (0.08,-0.16);
  }},
  pics/icoSegments/.style={code={
    \fill[black!16] (-0.20,0.10) rectangle (0.20,0.20);
    \fill[black!22] (-0.20,-0.05) rectangle (0.20,0.05);
    \fill[black!16] (-0.20,-0.20) rectangle (0.20,-0.10);
    \draw[black!50, line width=0.4pt] (-0.20,0.10) rectangle (0.20,0.20);
    \draw[black!50, line width=0.4pt] (-0.20,-0.05) rectangle (0.20,0.05);
    \draw[black!50, line width=0.4pt] (-0.20,-0.20) rectangle (0.20,-0.10);
  }}
}

% ----- Niveau 1 : objet fichier (MinIO) -----
\node[font=\sffamily\bfseries\tiny, text=black!45, anchor=east]
  at (0.15, 8.55) {MinIO};

\node[badge, fill=green!12] (tampon) at (2.55, 8.55) {};
\pic at (tampon.center) {icoDossier};
\node[below=8pt of tampon, align=center] {%
  \sffamily\scriptsize\bfseries Zone tampon\\[-0.5pt]
  \sffamily\tiny\color{black!50}\texttt{buffer/}};

\node[badge, fill=green!12] (chemin) at (10.35, 8.55) {};
\pic at (chemin.center) {icoDossier};
\node[below=8pt of chemin, align=center] {%
  \sffamily\scriptsize\bfseries Chemin calculé\\[-0.5pt]
  \sffamily\tiny\color{black!50}expéditeur / type / année / mois};

\draw[lien] (tampon) -- (chemin)
  node[midway, above=3pt, labfleche] {déplacement après analyse};
\fill[black!75] (tampon.east) circle (1.55pt);
\fill[black!75] (chemin.west) circle (1.55pt);

% Bifurcation vers PostgreSQL
\coordinate (jonction) at ($(tampon.east)!0.50!(chemin.west)$);
\fill[black!75] (jonction) circle (1.7pt);

% ----- Niveau 2 : description et classement (PostgreSQL) -----
\node[font=\sffamily\bfseries\tiny, text=black!45, anchor=east]
  at (0.15, 4.15) {PostgreSQL};

\node[badge, fill=green!10] (ref) at (1.85, 4.15) {};
\pic at (ref.center) {icoListe};
\node[below=8pt of ref, align=center] {%
  \sffamily\scriptsize\bfseries Référentiels\\[-0.5pt]
  \sffamily\tiny\color{black!50}types, organismes};

\node[badge, fill=blue!10] (doc) at (6.45, 4.15) {};
\pic at (doc.center) {icoPage};
\node[below=8pt of doc, align=center] {%
  \sffamily\scriptsize\bfseries Document\\[-0.5pt]
  \sffamily\tiny\color{black!50}champs, JSONB, empreinte};

\node[badge, fill=blue!10] (noeud) at (10.35, 4.15) {};
\pic at (noeud.center) {icoArbre};
\node[below=8pt of noeud, align=center] {%
  \sffamily\scriptsize\bfseries Nœud\\[-0.5pt]
  \sffamily\tiny\color{black!50}DRAFT / VALIDATED / ARCHIVED};

\node[badge, fill=orange!12] (droits) at (14.15, 4.15) {};
\pic at (droits.center) {icoClef};
\node[below=8pt of droits, align=center] {%
  \sffamily\scriptsize\bfseries Droits\\[-0.5pt]
  \sffamily\tiny\color{black!50}rôles, ACL};

\draw[lien] (jonction) -- (doc.north)
  node[pos=0.42, labfleche, right=4pt, align=left] {référence\\de stockage};
\fill[black!75] (doc.north) circle (1.55pt);

\draw[lien] (ref) -- (doc)
  node[midway, above=3pt, labfleche] {vocabulaire};
\draw[lien] (doc) -- (noeud)
  node[midway, above=3pt, labfleche] {classement};
\draw[lien] (noeud) -- (droits)
  node[midway, above=3pt, labfleche] {accès};
\fill[black!75] (ref.east) circle (1.55pt);
\fill[black!75] (doc.west) circle (1.55pt);
\fill[black!75] (doc.east) circle (1.55pt);
\fill[black!75] (noeud.west) circle (1.55pt);
\fill[black!75] (noeud.east) circle (1.55pt);
\fill[black!75] (droits.west) circle (1.55pt);

% ----- Niveau 3 : index de recherche -----
\node[badge, fill=violet!12] (seg) at (6.45, 0.55) {};
\pic at (seg.center) {icoSegments};
\node[below=8pt of seg, align=center] {%
  \sffamily\scriptsize\bfseries Segments\\[-0.5pt]
  \sffamily\tiny\color{black!50}texte + vecteur (pgvector)};

\draw[lien] (doc.south) -- (seg.north)
  node[pos=0.48, labfleche, right=4pt, align=left] {index de\\recherche};
\fill[black!75] (doc.south) circle (1.55pt);
\fill[black!75] (seg.north) circle (1.55pt);

\end{tikzpicture}%
}
\endgroup
\vspace{0.12em}
\footnotesize Le fichier reste un objet MinIO ; PostgreSQL conserve sa référence de stockage ainsi que la description, le classement, l'index de recherche et les informations d'accès.
